\chapter{Magnetorquer Controller}

As mentioned in chapters \ref{chap:modelling}, magnetorquers are the only mechanism for dumping momentum build up on a satellite with reaction wheels. Furthermore, the FNDO-BS controller derived in chapter \ref{chap:underactuated_controller} requires that the angular momentum be zero, i.e., $\vec{h}=0$. Finally due to the lack of authority of said controller on the underactuated axis near the equilibrium, small steady-state errors can build over time during manoeuvres on the dead axis.  In this chapter, a magnetorquer controller is presented which must address these issues. A simple allocation scheme that achieves both objectives, and demonstrates its performance in simulation.
%@TODO: reword this intro

\section{Torque Generation}
A simple proportional controller is used to generate a torque that is opposite to the angular momentum of the system. Reiterating equation \eqref{eq:magt_torque} given in chapter \ref{chap:modelling}, torque generated by the magnetorquers $\vec{\tau}_m$ is given by \cite{sidiAttitudeDynamicsKinematics1997}:

% is this body frame?
\begin{equation*}
    \vec{\tau}_m = \vec{m} \times \vec{B}
    \label{eq:magt_torque_2}
\end{equation*}

where $\vec{m}$ is the magnetic moment of the magnetorquers and $\vec{B}$ is the magnetic field strength at the satellite's location in body coordinates. Equation \eqref{eq:magt_torque} is not invertible, however with some simplifying assumptions a magnetic moment $m$ can be determined. Assuming that the commanded torque is always perpendicular to the local field - $\vec{\tau}_m \perp \vec{B}$ for every commanded dipole. This means the set of torques reachable at any instant is limited to the plane normal to the local field rather than the whole of $\mathbb{R}^3$. Writing $\hat{\vec{B}} = \vec{B}/\|\vec{B}\|$, the minimum-norm dipole realising a commanded torque $\vec{\tau}_c$ is \cite{sidiAttitudeDynamicsKinematics1997}

\begin{equation}
    \vec{m} = \frac{\vec{B} \times \vec{\tau}_c}{\|\vec{B}\|^2}
    \label{eq:magt_alloc}
\end{equation}

and the torque actually delivered is the orthogonal projection of the demand onto that plane,
\begin{equation}
    \vec{\tau}_m = \vec{m} \times \vec{B} = \left(\vec{I}_3 - \hat{\vec{B}}\hat{\vec{B}}^T\right)\vec{\tau}_c \triangleq \vec{P}_B\,\vec{\tau}_c
    \label{eq:magt_projection}
\end{equation}
 A torque demand parallel to the instantaneous field direction produces a commanded dipole that is zero. This is can be observed mathematically by noting that the eigenvalues of $\vec{P}_B$ are ($1, 1, 0$) which means it has rank 2.
\section{Momentum Dumping}
\label{subsec:null_space_parking}
The conventional momentum dumping control law sets the demand proportional to the total angular momentum, 
\begin{equation}
    \vec{\tau}_c = -k\,\vec{\Delta h}, \ k > 0
\end{equation}
where $\vec{\Delta h} = \vec{h} - \vec{h}_d$ is the difference between the current and desired angular momentum.

Substituting $\vec{\tau}_c$ into equation \eqref{eq:magt_projection}, only the component of $\vec{h}$ perpendicular to the field is removed \cite{sidiAttitudeDynamicsKinematics1997},
\begin{equation}
    \dot{\vec{h}} = -k\,\vec{P}_B\,\vec{\Delta h} = -k\left(\vec{h} - \hat{\vec{B}}(\hat{\vec{B}}^T\vec{\Delta h})\right)
    \label{eq:dump_continuous}
\end{equation}
Because the law is implemented at a sample period $\Delta t$ with a zero-order hold, the discrete update is
\begin{equation}
    \vec{h}_{k+1} = (1 - k\Delta t)\,\vec{h}_k + k\Delta t\,\hat{\vec{B}}_k(\hat{\vec{B}}_k^T\vec{h}_k)
    \label{eq:dump_discrete}
\end{equation}
Note that $k$ carries units of $\mathrm{s}^{-1}$ with $\vec{\tau}_c = -k\vec{h}$ the gain \emph{is} the commanded decay rate, not a dimensionless tuning constant. Equation \eqref{eq:dump_discrete} then exposes a failure mode that is easily mistaken for insufficient actuator authority. Setting $k\Delta t = 1$ gives the deadbeat case
\begin{equation}
    \vec{h}_{k+1} = \hat{\vec{B}}_k(\hat{\vec{B}}_k^T\vec{h}_k)
\end{equation}
After a single sample the entire surviving momentum lies entirely along $\hat{\vec{B}}_k$, that is, inside the null space of $\vec{P}_B$. On the next cycle the demand is parallel to the field, the commanded dipole is zero, and the controller stalls. The momentum has been driven into the one direction the actuator cannot address. Intuitively, a greedy law removes all the momentum it can observe, and what it leaves behind is precisely the part it cannot observe.

Escape is possible only because $\hat{\vec{B}}$ rotates in the body frame. For a near-circular orbit this rotation occurs at approximately twice the orbital rate, $\omega_B \approx 2n$, which for the 500\,km, $i \approx 90^\circ$ orbit used throughout this work is $\omega_B \approx 2.2\times10^{-3}\,\mathrm{rad/s}$. Over one sample the field direction moves by $\varepsilon = \omega_B\Delta t$, so the momentum left pointing at the previous field direction is no longer quite aligned with the new one and a fraction $\sin\varepsilon$ becomes visible again. Under deadbeat action all of it is removed, leaving
\begin{equation}
    \|\vec{h}_{k+1}\| = \|\vec{h}_k\|\cos\varepsilon \approx \|\vec{h}_k\|\left(1 - \tfrac{1}{2}\varepsilon^2\right)
    \label{eq:parking_leak}
\end{equation}
The leak is \emph{quadratic} in the sweep angle. For $\Delta t = 1\,\mathrm{s}$ this gives a decay rate of $2.4\times10^{-6}\,\mathrm{s}^{-1}$, a time constant of roughly $4.8$ days. A simulation of the law at $k = 1$, $\Delta t = 1\,\mathrm{s}$ reproduces this: the residual momentum settles at $0.00^\circ$ from the field direction and decays with a fitted time constant of $4.6$ days, effectively frozen on any timescale of operational interest.

Reducing the gain so that $k\Delta t \ll 1$ prevents the collapse. The momentum then never aligns with the field but settles at an equilibrium lag angle where the rate of removal balances the rate at which the field sweeps past, and the decay becomes linear rather than quadratic in the geometry,
\begin{equation}
    \dot{\|\vec{h}\|} \approx -k\,\|\vec{h}\|\sin^2\theta, \qquad \theta = \angle(\vec{h}, \hat{\vec{B}})
    \label{eq:dump_linear_decay}
\end{equation}
\autoref{tab:dump_gain_sweep} shows this over a sweep of $k$, with the spacecraft inertia $\vec{J} = \mathrm{diag}(2.5, 3.5, 4.5)\,\mathrm{kg\,m^2}$ and an initial momentum of $0.0866\,\mathrm{N\,m\,s}$. The lag angle column is the diagnostic: at high gain the residual sits at $0^\circ$ from the field and is unreachable, while at $k = 3\times10^{-3}$ it rides at roughly $66^\circ$ and remains fully actuatable.

\begin{table}[h]
\centering
\begin{tabular}{@{}l l l l l@{}}
\hline
$k$ [s$^{-1}$] & $k\Delta t$ & $\|\vec{h}\|$ @ 2 orbits & $\|\vec{h}\|$ @ 5 orbits & $\angle(\vec{h}, \hat{\vec{B}})$ \\
\hline
$1\times10^{0}$  & 1.000 & $6.10\times10^{-3}$ & $5.85\times10^{-3}$ & $0.00^\circ$ \\
$1\times10^{-1}$ & 0.100 & $2.55\times10^{-3}$ & $1.13\times10^{-3}$ & $1.62^\circ$ \\
$3\times10^{-2}$ & 0.030 & $1.20\times10^{-4}$ & $7.17\times10^{-6}$ & $5.77^\circ$ \\
$1\times10^{-2}$ & 0.010 & $3.94\times10^{-5}$ & $4.54\times10^{-9}$ & $17.62^\circ$ \\
$3\times10^{-3}$ & 0.003 & $4.39\times10^{-9}$ & $2.26\times10^{-20}$ & $65.65^\circ$ \\
$1\times10^{-3}$ & 0.001 & $2.86\times10^{-4}$ & $4.63\times10^{-8}$ & $66.99^\circ$ \\
$3\times10^{-4}$ & 0.000 & $1.44\times10^{-2}$ & $1.09\times10^{-3}$ & $28.62^\circ$ \\
\hline
\end{tabular}
\caption{Residual angular momentum after magnetic dumping, as a function of the dumping gain, at $\Delta t = 1\,\mathrm{s}$ and a dipole limit of $1000\,\mathrm{A\,m^2}$. Performance is non-monotonic in $k$: too high parks the residual in the null space, too low is simply slow.}
\label{tab:dump_gain_sweep}
\end{table}

The resulting design rule is that the dumping gain must be \emph{slower} than the rate at which the field sweeps through the body frame,
\begin{equation}
    k \lesssim \omega_B \approx 2n
    \label{eq:dump_gain_rule}
\end{equation}
which corresponds to a dumping time constant on the order of the orbital period. This is counter-intuitive relative to the usual instinct to raise a regulator gain when convergence is poor, and it is worth stating explicitly because the symptom of violating it is a commanded dipole that collapses towards zero, which presents as though the actuator were undersized.

The same sweep shows that the available dipole is not the limiting factor at all. At $k = 3\times10^{-3}$ the dumping converges identically for dipole limits of $1000$, $100$, $30$ and $10\,\mathrm{A\,m^2}$, and still converges at $1\,\mathrm{A\,m^2}$; the peak dipole actually demanded is under $10\,\mathrm{A\,m^2}$ in every case.

\subsection{Dipole Saturation}
Because equation \eqref{eq:magt_alloc} can demand an arbitrarily large dipole as $\|\vec{B}\|$ falls or the demand grows, the commanded moment must be limited to the physically available $m_{max}$. Saturating each body component independently, $m_i \leftarrow \mathrm{sat}(m_i, m_{max})$, rotates $\vec{m}$ away from the direction $\vec{B} \times \vec{\tau}_c$, so the delivered torque $\vec{m} \times \vec{B}$ is no longer parallel to the momentum being removed and can add momentum on individual axes. Scaling the vector instead,
\begin{equation}
    \vec{m} \leftarrow \vec{m}\,\min\!\left(1, \frac{m_{max}}{\|\vec{m}\|}\right)
    \label{eq:dipole_scaling}
\end{equation}
preserves the torque direction and reduces only its magnitude, which is the correct behaviour under authority limits.

% \subsection{Orbit-Averaged Controllability}
% Since $\vec{P}_B$ is rank deficient at every instant, dumping in all three axes is possible only in an averaged sense as the field sweeps through the body frame. The relevant quantity is the time average of the projector,
% \begin{equation}
%     \bar{\vec{P}}_B(T) = \frac{1}{T}\int_0^{T}\left(\vec{I}_3 - \hat{\vec{B}}(t)\hat{\vec{B}}(t)^T\right)dt
%     \label{eq:avg_projector}
% \end{equation}
% whose smallest eigenvalue measures the authority available in the least favourable direction over a window of length $T$. \autoref{tab:avg_projector} evaluates this for the 500\,km orbit using the tilted-dipole field model of \autoref{chap:modelling}.

% \begin{table}[h]
% \centering
% \begin{tabular}{@{}l l l@{}}
% \hline
% Window $T$ & Field arc swept & $\lambda_{\min}(\bar{\vec{P}}_B)$ \\
% \hline
% $300\,\mathrm{s}$   & $43.2^\circ$  & 0.047 \\
% $600\,\mathrm{s}$   & $78.0^\circ$  & 0.141 \\
% $1200\,\mathrm{s}$  & $136.8^\circ$ & 0.339 \\
% $0.25\,T_{orb}$     & $157.7^\circ$ & 0.402 \\
% $0.5\,T_{orb}$      & $176.2^\circ$ & 0.499 \\
% $1\,T_{orb}$        & $176.2^\circ$ & 0.500 \\
% \hline
% \end{tabular}
% \caption{Smallest eigenvalue of the orbit-averaged projector as a function of averaging window, at 500\,km and $i \approx 90^\circ$. Over a full orbit no direction is starved, but over short windows one direction is close to unreachable.}
% \label{tab:avg_projector}
% \end{table}

% Over a full orbit the average is well conditioned and no direction is starved, confirming that complete magnetic dumping is achievable in principle. Over a $300\,\mathrm{s}$ window, however, the field sweeps only $43^\circ$ and the weakest direction retains under $5\%$ of isotropic authority. Magnetic momentum management is therefore an orbit-timescale process, and a simulation or operational window shorter than roughly one fifth of an orbit cannot dump all three axes regardless of gain or dipole capacity. This has a direct bearing on the controllability assumption of \autoref{chap:underactuated_controller}: the requirement $\vec{h} = \vec{0}$ inherited from \cite{aeyelsStabilizationNonLinear} cannot be established quickly by magnetic means, so any controller whose stability argument presumes it must tolerate a non-zero $\vec{h}$ for a substantial fraction of an orbit.

\subsection{The Unactuated Axis Is the Binding Constraint}
\label{subsec:dead_axis_constraint}
The preceding analysis concerns $\|\vec{h}\|$ as a whole, but for the underactuated spacecraft the quantity that actually determines pointing performance is the single component of $\vec{h}$ along the unactuated axis. Reaction wheel torques are internal and cannot change $\vec{h}$ at all; they only redistribute it between body and wheels. With the $z$-axis wheel failed, $h_{wz} = \dot{h}_{wz} = 0$, so
\begin{equation}
    J_z\,\omega_z = h_z - h_{wz} = h_z
    \label{eq:dead_axis_rate}
\end{equation}
and the yaw rate is fixed entirely by the projection of the total momentum onto the current body $z$ axis. Consequently the yaw angle is the double integral of the dead-axis disturbance torque, with no restoring term of any kind,
\begin{equation}
    \psi(t) = \frac{1}{J_z}\int_0^t h_z(\tau)\,d\tau, \qquad \ddot{\psi} = \frac{T_{d,z}}{J_z}
    \label{eq:yaw_double_integral}
\end{equation}
% Closed-loop simulation confirms equation \eqref{eq:yaw_double_integral} exactly: over a $6000\,\mathrm{s}$ inertial-hold run the yaw reconstructed from $\frac{1}{J_z}\int h_z\,dt$ matches the logged attitude with an RMS mismatch of $0.00^\circ$, both peaking at $39.18^\circ$ at $t = 2207\,\mathrm{s}$.

% Two consequences follow. First, the dead axis is the most fragile despite carrying the \emph{smallest} disturbance. In the same run the dead-axis disturbance amplitude was $1.9\times10^{-6}\,\mathrm{N\,m}$, some $38$ times smaller than the pitch-axis disturbance of $7.2\times10^{-5}\,\mathrm{N\,m}$, yet it alone produced tens of degrees of error because the actuated axes regulate their disturbance while the unactuated axis integrates it twice.

% Second, the condition for the dead axis to be instantaneously quiescent is weaker than the $\vec{h} = \vec{0}$ normally quoted. Equation \eqref{eq:dead_axis_rate} gives $\omega_z = 0$ whenever
% \begin{equation}
%     \vec{h}\cdot\hat{\vec{e}}_z = h_{wz}
%     \label{eq:dead_axis_condition}
% \end{equation}
% which is a single scalar constraint rather than three. This is, however, a pointwise condition and not an invariant one, and the distinction is decisive. In the absence of external torque $\vec{h}$ is fixed in \emph{inertial} space while $\hat{\vec{e}}_z$ is fixed in the \emph{body}, so the projection $\vec{h}\cdot\hat{\vec{e}}_z$ changes whenever the spacecraft slews. For the manoeuvre of \autoref{subsec:nadafi_tracking_sims}, a $70^\circ$ eigenaxis rotation about an axis lying in the body $xy$-plane, the body $z$ axis itself sweeps through $70^\circ$, and momentum that began entirely in the actuated plane is tipped onto the dead axis by as much as $\|\vec{h}\|\sin 70^\circ = 0.94\|\vec{h}\|$. Momentum stored in the actuated plane is therefore latent rather than disposed of. An inertial hold is precisely the special case in which it stays latent, because $\hat{\vec{e}}_z$ does not move and a yaw drift is a rotation \emph{about} $\hat{\vec{e}}_z$; this is why the verification of equation \eqref{eq:yaw_double_integral} above, which was performed on an inertial hold, is consistent with equation \eqref{eq:dead_axis_condition} while a tracking manoeuvre need not be.

% A second mechanism defeats the same shortcut independently of any slew of $\hat{\vec{e}}_z$. Holding attitude against a secular disturbance with $\vec{h} \neq \vec{0}$ obliges the two surviving wheels to carry a standing bias $h_{wx}, h_{wy} \neq 0$, and the dead-axis row of equation \eqref{eq:sat_dynamics} then reads
% \begin{equation}
%     J_z\dot{\omega}_z = -(J_y - J_x)\,\omega_x\omega_y \;-\; \left(\omega_x h_{wy} - \omega_y h_{wx}\right) \;+\; T_{d,z}
%     \label{eq:dead_axis_gyroscopic}
% \end{equation}
% so any actuated rate is converted into dead-axis angular acceleration through the wheel bias. This is not an incidental coupling: the hypersurface of equation \eqref{eq:hypersurface} steers the unactuated axis precisely by commanding non-zero $\omega_{ex}$ and $\omega_{ey}$, so the manoeuvre the controller requires in order to function is also the manoeuvre that excites this term. A wheel bias is an active source on the one axis that has no actuator, not passive storage.

% Both mechanisms trace to the same step in the derivation. Assumption~1 of \cite{nadafiRobustBacksteppingAttitude2021} discards the gyroscopic terms $\vec{\omega}^\times\vec{J}\vec{\omega}$ and $\vec{\omega}^\times\vec{h}_w$ from the actuated rows of equation \eqref{eq:sat_dynamics}, and with $h_{wz} = \dot{h}_{wz} = 0$ every surviving entry of both terms in those two rows carries a factor of $\omega_z$. The cascade model of equation \eqref{eq:cascade_dyn} is therefore exact only while $\omega_z \equiv 0$, which by equation \eqref{eq:dead_axis_rate} demands $h_z \equiv 0$ along the \emph{whole} trajectory; and since $h_z(t) = \hat{\vec{e}}_z^T\vec{T}_{BI}(t)\,\vec{h}_I$ must then vanish for every attitude the trajectory visits, this recovers $\vec{h} = \vec{0}$ rather than the weaker $\vec{h} \perp \hat{\vec{e}}_z$ at a single instant. Equation \eqref{eq:dead_axis_condition} is thus a diagnostic for when the dead axis is momentarily quiescent, and a design target for the magnetorquer on an inertial hold, but it is not a relaxation of the assumption under which the control law is derived.


\subsection{Split Allocation}
\label{subsec:split_allocation}
The observations above motivate treating the magnetorquer not as a momentum dumper that incidentally keeps the unactuated axis clean, but as the dedicated actuator on that axis. Since the dead-axis dynamics are the double integrator of equation \eqref{eq:yaw_double_integral}, a proportional-derivative law is appropriate. Using the vector part of the body-frame attitude error, for which $q_{e,z} \approx -\psi/2$, and noting from equation \eqref{eq:dead_axis_rate} that $h_z$ is directly proportional to the yaw rate,
\begin{equation}
    T_{d,z} = -2\zeta\omega_n\,h_z + 2J_z\omega_n^2\,q_{e,z}
    \label{eq:dead_axis_pd}
\end{equation}
giving the closed loop $\ddot{\psi} + 2\zeta\omega_n\dot{\psi} + \omega_n^2\psi = T_{dist,z}/J_z$ and hence a steady-state error of
\begin{equation}
    \psi_{ss} \approx \frac{T_{dist,z}}{J_z\,\omega_n^2}
    \label{eq:yaw_ss_error}
\end{equation}
The remaining two axes retain the slow dumping law of \autoref{subsec:null_space_parking} applied to $\vec{h}$ with its dead-axis component removed. The bandwidth $\omega_n$ must lie well above the orbital rate so that the dead-axis disturbance is rejected, and well below $1/\Delta t$.

Combining the two demands requires care. Naively summing them and passing the result through equation \eqref{eq:magt_alloc} does \emph{not} decouple them, because the dumping demand on the actuated axes projects a spurious torque onto the dead axis. In simulation this leakage was of order $3\times10^{-5}\,\mathrm{N\,m}$, more than ten times the $1.9\times10^{-6}\,\mathrm{N\,m}$ disturbance the dead-axis loop is attempting to reject, so the yaw error settled where the PD command balanced the leakage rather than the disturbance, degrading the achieved pointing by roughly two orders of magnitude.

Exact decoupling follows from observing that the dead-axis torque is
\begin{equation}
    \vec{\tau}_m\cdot\hat{\vec{e}}_z = (\vec{m}\times\vec{B})\cdot\hat{\vec{e}}_z = \vec{m}\cdot(\vec{B}\times\hat{\vec{e}}_z) \triangleq \vec{m}\cdot\vec{u}
    \label{eq:dead_axis_direction}
\end{equation}
so that only the component of $\vec{m}$ along $\vec{u} = \vec{B}\times\hat{\vec{e}}_z$ produces any dead-axis torque, and every dipole perpendicular to $\vec{u}$ produces none. Splitting the commanded dipole on $\vec{u}$ therefore separates the two objectives,
\begin{align}
    \vec{m}_{\psi} &= \frac{T_{d,z}}{\|\vec{u}\|^2}\,\vec{u} \label{eq:m_yaw}\\
    \vec{m}_{dump} &= \frac{\vec{B}\times\vec{T}_{d,\perp}}{\|\vec{B}\|^2} - \frac{(\vec{B}\times\vec{T}_{d,\perp})\cdot\vec{u}}{\|\vec{B}\|^2\,\|\vec{u}\|^2}\,\vec{u} \label{eq:m_dump}\\
    \vec{m} &= \vec{m}_{\psi} + \vec{m}_{dump} \label{eq:m_total}
\end{align}
where $\vec{T}_{d,\perp}$ is the dumping demand on the actuated axes. Since $\vec{m}_{\psi} \parallel \vec{u}$ and $\vec{m}_{dump} \perp \vec{u}$ by construction, the two are orthogonal and saturation with the dead axis given priority has the closed form
\begin{equation}
    \vec{m}_{dump} \leftarrow \vec{m}_{dump}\min\!\left(1, \frac{\sqrt{m_{max}^2 - \|\vec{m}_{\psi}\|^2}}{\|\vec{m}_{dump}\|}\right)
    \label{eq:priority_saturation}
\end{equation}
with $\vec{m}_{\psi}$ itself scaled back only if it alone exceeds $m_{max}$. In the degenerate case $\vec{B} \parallel \hat{\vec{e}}_z$ the vector $\vec{u}$ vanishes and no dead-axis authority exists for that cycle; the yaw term is then simply suppressed. For a nadir-pointing spacecraft this occurs where the field is radial, that is, near the magnetic poles. In the orbit studied the minimum angle between $\vec{B}$ and the dead axis was $1.5^\circ$ and the available dead-axis torque never fell below a factor of $59$ above the disturbance, so these windows proved to be of no practical consequence.

\subsection{Simulation Results}
\autoref{tab:split_allocation_results} compares the conventional isotropic dumping law against the split allocation on an inertially-referenced hold, driven by the disturbance torque history recorded from the full closed-loop simulation. Because reaction wheels cannot change $\vec{h}$, the dead-axis model used for this comparison is exact given the disturbance history.

\begin{table}[h]
\centering
\begin{tabular}{@{}l l l@{}}
\hline
Configuration & Peak $|\psi|$ & Peak $\|\vec{m}\|$ \\
\hline
Isotropic dumping, $k = 3\times10^{-3}$          & $48.99^\circ$ & $2.05\,\mathrm{A\,m^2}$ \\
Split allocation, $\omega_n = 0.005\,\mathrm{rad/s}$ & $1.04^\circ$  & $4.20\,\mathrm{A\,m^2}$ \\
Split allocation, $\omega_n = 0.01\,\mathrm{rad/s}$  & $0.27^\circ$  & $4.19\,\mathrm{A\,m^2}$ \\
Split allocation, $\omega_n = 0.02\,\mathrm{rad/s}$  & $0.07^\circ$  & $4.20\,\mathrm{A\,m^2}$ \\
\hline
\end{tabular}
\caption{Unactuated-axis pointing error over a $6000\,\mathrm{s}$ inertial hold. The $\omega_n = 0.01$ result agrees with the $0.24^\circ$ predicted by equation \eqref{eq:yaw_ss_error}.}
\label{tab:split_allocation_results}
\end{table}

The split allocation improves unactuated-axis pointing by more than two orders of magnitude, and the measured error agrees with the analytical prediction of equation \eqref{eq:yaw_ss_error}. The peak dipole demand is approximately $4\,\mathrm{A\,m^2}$ in all cases and the results are unchanged for $m_{max}$ anywhere from $100$ down to $5\,\mathrm{A\,m^2}$, confirming that the improvement comes from the allocation rather than from additional actuator capacity.

A separate result emerged when the same allocation was exercised during a commanded yaw manoeuvre rather than a hold. Because the wheels cannot produce dead-axis torque, the box manoeuvre steers yaw by pumping $h_z$ through the gyroscopic term $-(\vec\omega\times\vec{h})_z$, which is precisely the quantity the rate feedback in equation \eqref{eq:dead_axis_pd} damps. It was therefore anticipated that the dead-axis loop would oppose the manoeuvre and that detuning it during the manoeuvre would improve tracking. \autoref{tab:manoeuvre_gain} shows the opposite.

\begin{table}[h]
\centering
\begin{tabular}{@{}l l l@{}}
\hline
Dead-axis gain during manoeuvre & $q_{e,z}$ at $t = 240\,\mathrm{s}$ & Yaw error \\
\hline
$\omega_n = 0.015$ (unchanged)      & 0.1753  & $+0.19^\circ$  \\
$\omega_n = 0.002$ (detuned)        & 0.0695  & $-11.95^\circ$ \\
$\omega_n = 0$ (suppressed)         & $-0.2889$ & $-55.10^\circ$ \\
$\omega_n = 0$ throughout           & $-0.3662$ & $-65.34^\circ$ \\
\hline
\end{tabular}
\caption{Yaw tracking through a box manoeuvre at $m_{max} = 15\,\mathrm{A\,m^2}$, target $q_{e,z} = 0.1736$. Tracking improves monotonically with dead-axis magnetic authority.}
\label{tab:manoeuvre_gain}
\end{table}

Tracking improves monotonically with dead-axis magnetic authority: with the loop suppressed the manoeuvre diverges to $-65^\circ$, in the wrong direction entirely, while with it fully active the error settles to $+0.19^\circ$. The magnetorquer is not competing with the wheels for yaw control but supplying the majority of it. The in-plane momentum available to the manoeuvre was within $15\%$ across these cases, so the difference is attributable to the dead-axis loop rather than to the dumping law depriving the manoeuvre of the momentum it requires.

This also clarifies an earlier observation that the manoeuvre overshot its target by $168\%$ and then recovered only slowly. That behaviour occurred at an unrealistic $m_{max} = 1000\,\mathrm{A\,m^2}$; at a representative $15\,\mathrm{A\,m^2}$ the same configuration tracks to within $0.2^\circ$. The overshoot was a symptom of excessive dead-axis authority overdriving the loop, not of contention between the two actuators.

\subsection{Limitations}
Three limitations of the above should be noted. The gain sweep of \autoref{tab:dump_gain_sweep} and the hold results of \autoref{tab:split_allocation_results} were obtained from a reduced-order replay driven by logged disturbance and field histories; this is exact for the unactuated axis, since reaction wheels cannot alter $\vec{h}$, but does not capture the reaction of the backstepping controller and the fault observer to the magnetic torque appearing on the actuated axes. The manoeuvre results of \autoref{tab:manoeuvre_gain} are from full closed-loop simulation but represent a single trajectory at one orbital phase, so the quoted figures should be read as indicative rather than as general performance bounds. Finally, the loop of equation \eqref{eq:dead_axis_pd} is normalised in gain with respect to inertia but not in authority: the dipole required scales as $J_z\omega_n^2\psi/(\|\vec{B}\|\sin\angle(\vec{B},\hat{\vec{e}}_z))$, so a spacecraft of substantially larger moment of inertia will require proportionally more dipole to sustain the same bandwidth, and the $15\,\mathrm{A\,m^2}$ used here may not be sufficient.